% Options for packages loaded elsewhere
\PassOptionsToPackage{unicode}{hyperref}
\PassOptionsToPackage{hyphens}{url}
%
\documentclass[
]{article}
\usepackage{amsmath,amssymb}
\usepackage{iftex}
\ifPDFTeX
  \usepackage[T1]{fontenc}
  \usepackage[utf8]{inputenc}
  \usepackage{textcomp} % provide euro and other symbols
\else % if luatex or xetex
  \usepackage{unicode-math} % this also loads fontspec
  \defaultfontfeatures{Scale=MatchLowercase}
  \defaultfontfeatures[\rmfamily]{Ligatures=TeX,Scale=1}
\fi
\usepackage{lmodern}
\ifPDFTeX\else
  % xetex/luatex font selection
\fi
% Use upquote if available, for straight quotes in verbatim environments
\IfFileExists{upquote.sty}{\usepackage{upquote}}{}
\IfFileExists{microtype.sty}{% use microtype if available
  \usepackage[]{microtype}
  \UseMicrotypeSet[protrusion]{basicmath} % disable protrusion for tt fonts
}{}
\makeatletter
\@ifundefined{KOMAClassName}{% if non-KOMA class
  \IfFileExists{parskip.sty}{%
    \usepackage{parskip}
  }{% else
    \setlength{\parindent}{0pt}
    \setlength{\parskip}{6pt plus 2pt minus 1pt}}
}{% if KOMA class
  \KOMAoptions{parskip=half}}
\makeatother
\usepackage{xcolor}
\usepackage[margin=1in]{geometry}
\usepackage{graphicx}
\makeatletter
\newsavebox\pandoc@box
\newcommand*\pandocbounded[1]{% scales image to fit in text height/width
  \sbox\pandoc@box{#1}%
  \Gscale@div\@tempa{\textheight}{\dimexpr\ht\pandoc@box+\dp\pandoc@box\relax}%
  \Gscale@div\@tempb{\linewidth}{\wd\pandoc@box}%
  \ifdim\@tempb\p@<\@tempa\p@\let\@tempa\@tempb\fi% select the smaller of both
  \ifdim\@tempa\p@<\p@\scalebox{\@tempa}{\usebox\pandoc@box}%
  \else\usebox{\pandoc@box}%
  \fi%
}
% Set default figure placement to htbp
\def\fps@figure{htbp}
\makeatother
\setlength{\emergencystretch}{3em} % prevent overfull lines
\providecommand{\tightlist}{%
  \setlength{\itemsep}{0pt}\setlength{\parskip}{0pt}}
\setcounter{secnumdepth}{-\maxdimen} % remove section numbering
\usepackage{graphicx}
\usepackage{float}
\usepackage{bookmark}
\IfFileExists{xurl.sty}{\usepackage{xurl}}{} % add URL line breaks if available
\urlstyle{same}
\hypersetup{
  hidelinks,
  pdfcreator={LaTeX via pandoc}}

\author{}
\date{\vspace{-2.5em}}

\begin{document}

\section{Question}\label{question}

En el diagrama se muestran los valores proporcionales de participantes
en un entrenamiento técnico, dependiendo del género y la edad.

\includegraphics[width=7cm,height=\textheight,keepaspectratio]{tabla_contingencia_corregida.png}

Por ejemplo, el 26.8\% de los participantes son participantes de género
femenino mayores de 17 años. Según el diagrama, ¿cuál es la probabilidad
de que al escoger una persona al azar pertenezca al grupo de los menores
de 17 años, si ya se sabe que hace parte del grupo de participantes
masculinos?

\subsection{Answerlist}\label{answerlist}

\begin{itemize}
\tightlist
\item
  0.196/0.464
\item
  0.361/0.464
\item
  0.361/0.629
\item
  0.175/0.536
\end{itemize}

\section{Solution}\label{solution}

Para resolver este problema de \textbf{probabilidad condicional},
necesitamos aplicar correctamente la fórmula de probabilidad condicional
y trabajar con la información de la tabla de contingencia.

\subsubsection{Paso 1: Identificar el tipo de problema y la
fórmula}\label{paso-1-identificar-el-tipo-de-problema-y-la-fuxf3rmula}

Este es un problema de \textbf{probabilidad condicional}, donde
buscamos:

\[P(\text{menores de 17 años} | \text{participantes masculinos})\]

La fórmula de probabilidad condicional es:
\[P(A|B) = \frac{P(A \cap B)}{P(B)}\]

Donde:

\begin{itemize}
\tightlist
\item
  \(A\) = evento de interés: menores de 17 años
\item
  \(B\) = condición dada: participantes masculinos
\item
  \(P(A \cap B)\) = probabilidad de que ocurran ambos eventos
  simultáneamente
\item
  \(P(B)\) = probabilidad marginal de la condición
\end{itemize}

\subsubsection{Paso 2: Extraer información de la tabla de
contingencia}\label{paso-2-extraer-informaciuxf3n-de-la-tabla-de-contingencia}

De la tabla podemos obtener las siguientes \textbf{probabilidades
conjuntas}:

\begin{itemize}
\tightlist
\item
  P(menores de 17 años \(\cap\) participantes masculinos) = 0.175
\item
  P(menores de 17 años \(\cap\) participantes de género femenino) =
  0.196
\item
  P(mayores de 17 años \(\cap\) participantes masculinos) = 0.361
\item
  P(mayores de 17 años \(\cap\) participantes de género femenino) =
  0.268
\end{itemize}

\subsubsection{Paso 3: Calcular probabilidades
marginales}\label{paso-3-calcular-probabilidades-marginales}

Las \textbf{probabilidades marginales} se obtienen sumando las
probabilidades conjuntas:

\textbf{Por género:}

\begin{itemize}
\tightlist
\item
  P(participantes masculinos) = 0.175 + 0.361 = 0.536
\item
  P(participantes de género femenino) = 0.196 + 0.268 = 0.464
\end{itemize}

\textbf{Por edad:}

\begin{itemize}
\tightlist
\item
  P(menores de 17 años) = 0.175 + 0.196 = 0.371
\item
  P(mayores de 17 años) = 0.361 + 0.268 = 0.629
\end{itemize}

\subsubsection{Paso 4: Aplicar la fórmula para nuestro problema
específico}\label{paso-4-aplicar-la-fuxf3rmula-para-nuestro-problema-especuxedfico}

Para nuestro problema:

\[P(\text{menores de 17 años} | \text{participantes masculinos}) = \frac{P(\text{menores de 17 años} \cap \text{participantes masculinos})}{P(\text{participantes masculinos})}\]

Sustituyendo los valores de la tabla:
\[P(\text{menores de 17 años} | \text{participantes masculinos}) = \frac{0.175}{0.536}\]

\subsubsection{Paso 5: Verificación y análisis de
distractores}\label{paso-5-verificaciuxf3n-y-anuxe1lisis-de-distractores}

\textbf{Respuesta correcta:} 0.175/0.536 = 0.326

\textbf{Análisis de errores conceptuales comunes (distractores):}

\begin{enumerate}
\def\labelenumi{\arabic{enumi}.}
\tightlist
\item
  \textbf{0.361/0.629}: Error de \textbf{intercambio conceptual} -
  confundir la dirección de la probabilidad condicional o usar
  probabilidades marginales incorrectas
\item
  \textbf{0.196/0.464}: Error de \textbf{aplicación incorrecta de la
  fórmula} - usar probabilidades conjuntas o marginales en posiciones
  incorrectas de la fracción
\item
  \textbf{0.361/0.464}: Error de \textbf{cálculo sutil} - aplicar
  operaciones incorrectas como complementos o usar valores de celdas
  adyacentes de la tabla
\end{enumerate}

\subsubsection{Conclusión}\label{conclusiuxf3n}

Por lo tanto, la probabilidad de que una persona pertenezca al grupo de
menores de 17 años, dado que es participantes masculinos, es
\textbf{0.175/0.536}.

Esta respuesta tiene sentido porque:

\begin{itemize}
\tightlist
\item
  El resultado está entre 0 y 1 (como toda probabilidad)
\item
  El numerador representa la intersección de los dos eventos
\item
  El denominador representa la probabilidad de la condición dada
\item
  La interpretación es coherente con el contexto del problema
\end{itemize}

\subsection{Answerlist}\label{answerlist-1}

\begin{itemize}
\tightlist
\item
  Falso
\item
  Falso
\item
  Falso
\item
  Verdadero
\end{itemize}

\section{Meta-information}\label{meta-information}

exname: probabilidad\_condicional\_tabla\_contingencia\_corregido
extype: schoice exsolution: 0001 exshuffle: TRUE exsection:
Probabilidad\textbar Probabilidad condicional\textbar Tablas de
contingencia\textbar Razonamiento matemático exextra{[}Type{]}: Cálculo
exextra{[}Program{]}: R exextra{[}Language{]}: es exextra{[}Level{]}: 3

\end{document}
